\section{Introduction}

Large language models are increasingly deployed in high-stakes applications where incorrect predictions carry significant costs --- medical diagnosis, legal advice, and autonomous systems. Selective prediction allows models to abstain when uncertainty is high, offering a principled approach to improving reliability without retraining. Current methods rely on maximum probability thresholding, which captures only the mode of a probability distribution and may miss multi-modal uncertainty patterns. We investigated whether entropy-based uncertainty quantification captures additional signal beyond max-probability, but discovered we could not test this hypothesis --- not because our method failed, but because model substitution (GPT-2 instead of Llama-2-7B) invalidated the statistical test itself.

Our work identifies model capacity as a binary gate for experiment validity in selective prediction research. Below a capacity threshold, correlation-based validation experiments cannot run because models produce zero-variance outcomes; above it, hypothesis testing becomes possible. This is distinct from the well-known relationship between model size and performance magnitude --- we show that capacity gates the validity of statistical tests, not just their power or precision.

The problem appears straightforward on the surface: practitioners need reliable uncertainty estimates for selective prediction without retraining or ensembles, and turn to single-forward-pass signals extracted from token probability distributions. Max-probability thresholding has emerged as the standard baseline, where predictions below a confidence threshold are rejected. However, this approach captures only the mode of the distribution --- the single most likely token --- and ignores the full distributional shape. Information theory suggests that entropy, which quantifies distribution flatness, should capture multi-modal uncertainty that max-probability misses.

The deeper problem emerges when we attempt to validate this intuition experimentally. Entropy and max-probability both derive from the same distribution, making them correlated by construction. The key question is whether entropy provides additional signal beyond max-probability in cases where the two metrics disagree --- specifically, when max-probability is high but entropy is also high, indicating a peaked distribution with significant secondary modes. Testing this requires experiments where models produce measurable variance in prediction correctness. If all predictions are either correct or incorrect, correlation-based statistical tests become undefined --- not weak, but mathematically impossible to compute.

This experimental prerequisite has not been explicitly discussed in selective prediction literature. Researchers typically assume any model can test uncertainty estimation methods, viewing model choice as affecting performance magnitude (better models have higher accuracy), not the validity of statistical tests themselves. Our work reveals that model capacity is not just a performance variable but an infrastructural prerequisite for correlation-based validation: below a threshold, the experiment cannot run; above it, hypothesis testing becomes possible.

We discovered this dependency when implementing our validation experiment for entropy-based selective prediction on TriviaQA, a factual question-answering benchmark. The experimental design specified Llama-2-7B (7 billion parameters), which achieves over 10\% accuracy on TriviaQA and should produce sufficient variance in correctness for correlation testing. However, implementation substituted GPT-2 (117 million parameters), a model 60$\times$ smaller and not designed for factual knowledge recall. GPT-2 produced 0\% exact-match accuracy on 500 TriviaQA examples --- every prediction was incorrect, creating zero variance in the correctness variable. Spearman correlation between entropy and correctness returned NaN (not a number), not a negative or non-significant value, because correlation measures co-variation and requires both variables to vary.

Critically, our infrastructure validation succeeded. We achieved 100\% extraction rate for entropy values (500/500 predictions yielded valid entropy measurements), confirming that token probability distributions are fully accessible from frozen model forward passes. We also confirmed that high max-probability, high-entropy disagreement cases exist in non-trivial proportions: 8.20\% of predictions fell into the Q3 quadrant (high max-prob, high entropy), exceeding our 5\% threshold. These results validate the technical infrastructure and statistical framework independent of the model's predictive performance. The quadrant analysis framework is sound; the extraction method works. Only the hypothesis test itself --- whether entropy correlates with correctness --- remains untested due to invalid test conditions.

This creates an unusual scientific contribution: we successfully validated infrastructure while failing to test the hypothesis, and in doing so, revealed an experimental design requirement not explicitly documented in prior work. We hypothesize that models below approximately 7 billion parameters invalidate correlation-based validation experiments on factual question-answering tasks because small models cannot provide the variance in correctness that statistical tests require (based on \citet{Roberts2020} showing knowledge capacity thresholds at 10B+ parameters), though we tested only GPT-2 (117M, 0\% accuracy) and have not empirically confirmed the threshold across intermediate scales. This is not a hypothesis refutation --- our hypothesis about entropy's superiority remains untested, awaiting proper experimental conditions. Rather, it is a methodological insight: selective prediction experiments have a dependency on model capacity that gates the validity of statistical analyses, not just their power or precision.

Our contributions are threefold. First, we provide empirical validation of the infrastructure for entropy-based selective prediction: token probability distributions are fully extractable (100\% success rate), and disagreement patterns between entropy and max-probability exist in practice (8.20\% of predictions). Second, we identify model capacity as a binary gate for experiment validity in selective prediction research: below threshold, correlation tests are undefined; above threshold, hypothesis testing becomes possible. This is a methodological contribution applicable beyond our specific hypothesis. Third, we establish a framework for distinguishing infrastructure validation from hypothesis validation, demonstrating that undefined test results reflect invalid conditions rather than substantive refutations.

The remainder of this paper is organized as follows. Section 2 reviews related work on selective prediction, uncertainty quantification, and entropy-based methods, positioning our methodological contribution within the literature. Section 3 describes our experimental design, including the entropy extraction method, quadrant analysis framework, and the unintended model substitution that exposed the capacity dependency. Section 4 presents our results: successful infrastructure validation alongside the correlation test failure due to zero-variance correctness. Section 5 discusses the implications of our findings for future selective prediction experiments and establishes guidelines for minimum model performance thresholds. Section 6 concludes with directions for future work, including re-testing our hypothesis under valid conditions and extending capacity threshold analysis to other task types.
